\begin{itemize}
\item \textbf{Diamagnetismus:} bez pole mají atomy nulový výsledný magnetický moment. Pole vyvolá malý moment opačného směru, takže magnetická indukce je menší než ve vakuu; $\chi_m<0$, $\mu_r<1$. Příklady: Cu, Ag, voda.
\item \textbf{Paramagnetismus:} atomy mají nenulové momenty, které jsou bez pole neuspořádané. V poli se částečně orientují souhlasně, takže indukce vzroste; $\chi_m>0$, $\mu_r>1$. Příklady: Al, Pt, O$_2$.
\item \textbf{Feromagnetismus:} existují spontánně zmagnetované domény. Vnější pole mění jejich velikost a orientaci, vzniká silná magnetizace s nasycením a hysterezí; po odstranění pole může zůstat remanence. Příklady: Fe, Co, Ni.
\end{itemize}
